\chapter{LITERATURE REVIEW}

\section{Traditional Air Quality Forecasting}
Traditional air quality forecasting methods include statistical time-series models, regression models, and interpolation techniques. These methods are useful for simple forecasting problems, but they may not effectively capture complex nonlinear and spatio-temporal pollutant patterns. Since air pollution is influenced by weather, wind, traffic, and regional interactions, station-wise forecasting alone may not be sufficient for accurate regional prediction.

Baseline models remain important in this project because they provide a clear reference point for evaluating advanced graph learning models. Persistence models, linear regression, random forest, \gls{svm}, \gls{lstm}, and \gls{gru} can be used to compare whether the proposed wind-aware \gls{stgnn} provides meaningful improvement.

\section{Deep Learning for Time-Series Forecasting}
Deep learning models such as \gls{lstm} and \gls{gru} are widely used for time-series forecasting because they can learn temporal dependencies from historical data. In air quality forecasting, these models can learn how pollutant concentration changes over time. However, normal \gls{lstm} and \gls{gru} models usually focus on temporal patterns and may not fully model spatial relationships between different monitoring stations.

Spatio-temporal forecasting research combines temporal sequence models with spatial models to address this limitation. For example, recent PM$_{2.5}$ forecasting work uses recurrent units with graph neural components to capture both time evolution and spatial diffusion across monitoring locations \cite{pandey2024spatialdiffusion}.

\section{Graph Neural Networks for Air Quality Forecasting}
\glspl{gnn} can model relationships between monitoring stations by representing the monitoring network as a graph. In this graph, stations are nodes and spatial relationships are edges. Spatio-temporal \gls{gnn} models combine graph learning with temporal models to capture both spatial and temporal dependencies. This makes them suitable for air quality forecasting because pollutant levels at one station can be affected by neighboring stations.

PM2.5-GNN is an example of a domain-knowledge enhanced graph neural model for PM$_{2.5}$ forecasting \cite{wang2020pm25gnn}. E-STGCN further shows that spatio-temporal graph convolutional approaches can be adapted for air pollution forecasting and can address extreme pollutant behavior in polluted cities \cite{panja2024estgcn}. These studies motivate the proposed project's use of graph-based spatial learning with temporal forecasting.

\section{Wind-Aware Dynamic Graph Modeling}
Many graph-based forecasting models use fixed edges based on distance or historical correlation. However, pollutant dispersion can change depending on wind speed and wind direction. A wind-aware graph can update edge weights dynamically based on whether wind flows from one station toward another. This can make the graph structure more realistic and physically meaningful.

In the proposed system, edge weights will combine distance, historical pollutant correlation, wind speed, and wind direction alignment. This approach is expected to better represent directional pollutant transport than a purely distance-based graph.

\section{Explainable AI in Environmental Forecasting}
Deep learning models are often considered black boxes. For environmental decision-support systems, explainability is important because users need to understand why a model predicts high pollution in a certain area. In this project, attention weights, feature importance, and neighboring station influence will be analyzed to provide interpretable explanations.

GNNExplainer and related explanation methods show how graph model predictions can be interpreted by identifying important graph structures and input features \cite{ying2019gnnexplainer}. For the proposed system, explanation outputs will be compared with wind direction, distance, weather features, and pollutant history to check whether the learned relationships are physically reasonable.

\section{Research Gap}
Existing air quality forecasting methods often focus either on temporal prediction or fixed spatial relationships. A practical gap remains in combining wind-aware dynamic graph construction, explainable prediction, and scenario-based mitigation analysis in one system. This project addresses the gap by integrating public air quality data, meteorological data, spatio-temporal graph learning, explanation, and an interactive dashboard.

\begin{table}[H]
    \centering
    \small
    \caption{Comparison of Related Approaches and Proposed System}
    \label{tab:related_air_quality}
    \renewcommand{\arraystretch}{1.2}
    \setlength{\tabcolsep}{4pt}

    \begin{tabularx}{0.94\linewidth}{|
        >{\raggedright\arraybackslash}p{2.8cm}|
        >{\raggedright\arraybackslash}p{4.4cm}|
        >{\raggedright\arraybackslash}X|}
        \hline
        \textbf{Approach} & \textbf{Primary Focus} & \textbf{Limitation / Difference from Proposed System} \\
        \hline

        Station-wise regression &
        Forecasts pollutant concentration for each monitoring station independently. &
        Does not capture spatial pollutant movement or influence from neighboring stations. \\
        \hline

        LSTM/GRU-based forecasting &
        Learns temporal patterns from historical pollutant and meteorological data. &
        Captures time dependency but may ignore station-to-station spatial influence. \\
        \hline

        Fixed-graph ST-GNN &
        Models spatial and temporal dependencies using a predefined graph structure. &
        Graph edges are usually static and may not adapt to changing wind speed and wind direction. \\
        \hline

        Wind-aware ST-GNN &
        Uses distance, wind direction, wind speed, and pollutant correlation to model dynamic station influence. &
        Provides more physically meaningful edge weighting for pollutant transport modeling. \\
        \hline

        Proposed system &
        Combines wind-aware forecasting, explainability, risk visualization, and what-if scenario analysis. &
        Extends forecasting results into an interpretable decision-support dashboard for air quality analysis. \\
        \hline

    \end{tabularx}
\end{table}
